

\documentclass[a4paper]{article} %

\date{} %

% Please, do not change any of the above lines

\usepackage{amssymb,amsmath,amsthm} 
% add other LaTeX packages you need here.

\begin{document}
	\setcounter{page}{1} % Please, do not delete this line
	
	% Your title
	\title{Old but NU: Finding Large Monoidal Intervals}
	
	% Authors
	\author{Tamás Ágoston, Miklós Dormán \\ %
		University of Szeged, Szeged, Hungary\\ \\ % Affiliation 1
		%Author 4 \\ % If any other author with different Affilation
		%University of ZZWW\\ \\ % Affiliation 2 (if needed)
		% New author \\
		% New affiliation \\
		% Add authors and affiliation as needed
		\texttt{agostont@server.math.u-szeged.hu} % Only one corresponding e-mail, please
	}%
	
	\maketitle
	
	% The abstract
	
	Let $A$ be a finite set and $M$ a transformation monoid on $A$.  For $n\geq 1$, let $\mathcal{O}_A^{(n)}$ denote the set of all $n$-ary operations on $A$. The set \[\{f\in \mathcal{O}_A^{(n)}\colon f(\mu_1(x),\dots,\mu_n(x))\in M \text{ for all } \mu_1,\dots,\mu_n \in M\} \] is called the stabilizer of $M$. The lattice of all clones on a $A$ is naturally partitioned into finitely many intervals, each having as its least element the clone generated by a transformation monoid and  as its greatest element the corresponding stabilizer. The problem of classifying transformation monoids according to the cardinalities of the corresponding monoidal intervals was posed by Á.~Szendrei in~\cite{TamasAgoston_AS}. 
	
	Combining the results of Krokhin \cite{AAK95} with computer-assisted computations, we show that near-unanimity relations give rise to monoidal intervals of cardinality $\lvert \mathcal{P}(\mathbb{N})\rvert$. 
	
	%This is a sample abstract for SSAOS 2023.
	%
	%Please, edit it and then upload it as a \LaTeX~file at the link included in the
	%email we sent you.
	%
	%The abstract should not exceed two pages, including the bibliography.
	
	\begin{thebibliography}{99}
		%\bibitem{mundici1986}
		%D.~Mundici.
		%\newblock \emph{Interpretation of {A}{F} ${C}^*$-algebras in {L}ukasziewicz sentential calculus,}
		%\newblock J. Functional Analysis \textbf{65} 15--53 (1986)
		%\bibitem{gratzer1998}
		%G.~Gr{\" a}tzer.
		%\newblock \emph{General Lattice Theory,}
		%\newblock Birkh{\" a}user, second edition, 1998.
		%
		\bibitem{TamasAgoston_AS}
		\textsc{Á. Szendrei}, \textit{Clones in universal algebra}, Séminare de mathématiques supérieures, Vol. 99, Les presses de l'Université de Montréal, Montréal, 1986.
		\bibitem{AAK95} 
		Krokhin, A. A.: 
		Monoidal intervals in lattices of clones, 
		Algebra and Logic, {\bf 34}, 155–168 (1995)
		
	\end{thebibliography}
	
\end{document}


